
Post-training of large language models for code generation has converged on two dominant paradigms: execution-feedback reinforcement learning, exemplified by GRPO \citep{Shaoetal2024}, and direct preference optimization \citep{Rafailovetal2023}. Both methods are routinely evaluated by reporting pass@1 on HumanEval+ \citep{Liuetal2023} or MBPP+ \citep{Liuetal2023}. This evaluation practice answers the outcome question — which method produces more correct code — but it does not answer a structurally distinct question: does a given alignment method cause the policy to move in qualitatively different ways relative to the base model?

Specifically, two scenarios produce higher pass@1 with different implications for generalization and interpretability. In the first scenario, the aligned policy reallocates probability mass toward control-flow and data-flow AST transformations — the structural elements that determine functional correctness. In the second scenario, the policy becomes more decisive about existing surface patterns without changing the underlying distribution over program structures. Pass@1 does not distinguish these two scenarios.

This distinction is the motivating problem for the present work. The central quantity of interest is \textit{structural efficiency of policy movement}: the amount of semantically meaningful code transformation (restricted to control-flow and data-flow AST nodes) per unit KL divergence from the base model. A method with high structural efficiency spends its policy change budget — measured in KL units — on transformations that affect program behavior. A method with low structural efficiency expends KL budget on surface-level changes (variable names, whitespace, comment style) that do not affect execution.

The tools needed to measure this quantity exist in isolation. The FA-AST taxonomy \citep{Wangetal2020} classifies Python AST nodes into control-flow, data-flow, and surface categories. The ZSS algorithm [Zhang and Shasha, 1989] computes minimum-cost tree edit distance and has been validated against established code similarity metrics \citep{Dingetal2024}. KL divergence from the base model is logged during standard RL fine-tuning. The contribution of this paper is to compose these tools into a unified, reproducible framework for structural efficiency measurement, to validate the framework end-to-end on real model checkpoints, and to report what was found — and what was not found — in a preliminary application.

The paper makes the following contributions. First, it formally defines structural efficiency as semantic AST edit distance per unit KL divergence and defines the Semantic Edit Proportion (SEP) as a complementary quantity capturing the fraction of edits that target semantically relevant nodes. Second, it implements and validates an end-to-end measurement pipeline comprising five modules: FA-AST node classification, ZSS semantic edit distance, KL-matched checkpoint selection, SEP analysis, and bootstrap statistical testing. Third, it reports a preliminary empirical observation that GRPO and DPO produce nearly identical SEP values (≈0.237 for both) in a run that was subsequently found to be severely underpowered due to checkpoint aliasing. Fourth, it documents checkpoint aliasing as a confound — in which a nominally 27-pair analysis collapsed to an effective sample size of approximately 2 due to insufficient checkpoint diversity — and provides a one-assertion safeguard that prevents this failure mode.
